\documentclass{article}
\usepackage{anyfontsize}
\usepackage[UTF8]{ctex} % 如果需要支持中文
\usepackage{fontspec} 
% \setCJKmainfont{SourceHanSansCN-Light}[
%     BoldFont=SourceHanSansCN-Medium
% ]

\usepackage[a4paper, left=0.1in, right=0.1in, top=0.05in, bottom=0.05in]{geometry} % 设置四边页边距为0.1英寸{geometry} % 设置页边距为0.5英寸
\usepackage{enumitem} % 用于自定义列表
\usepackage{amsmath}  % 数学公式支持
\usepackage{tikz}
\usetikzlibrary{arrows.meta, positioning}
\setlength{\parindent}{0pt} % 不缩进
\usepackage{etoolbox} % 用于列表操作
%调整类代码
\usepackage{comment}
\usepackage{tocloft} % 用于定制目录 样式
\usepackage{hyperref} %不需要目录盒子注释这
\usepackage{ifthen}
\usepackage{tikz}
\usetikzlibrary{arrows.meta}
\usepackage{titletoc}
\usepackage{graphicx}
\usepackage{float}

\usepackage{amssymb}  % 提供数学符号，包括\therefore
%目录设定
%快捷键 CMD + / 注释
%  \renewcommand{\cftdot}{} % 隐藏目录中的页码
%  \cftpagenumbersoff{section}
%  \cftpagenumbersoff{subsection}
%  \makeatletter
%  \renewcommand{\@pnumwidth}{0em}  % 设置页码宽度为0
%  \renewcommand{\@tocrmarg}{0em}   % 设置右边距为0
%  \makeatother
 

\usepackage{environ}

\newif\ifshowcontent  % 定义一个条件判断变量
\showcontenttrue    % 设置为 false，隐藏所有 question 环境的内容

\newcounter{question} % 题目计数器
\setcounter{question}{0}
\NewEnviron{question}{%
    \ifshowcontent
        \par\stepcounter{question}\textbf{\thequestion.} \BODY\par % 如果显示内容，则输出
    \fi
}

\NewEnviron{choices}{%
    \ifshowcontent
        \begin{enumerate}[label=\Alph*., leftmargin=*]
            \BODY
        \end{enumerate}\par
    \fi
}

\NewEnviron{correctanswer}{%
    \ifshowcontent
        \textbf{答案:}\BODY\par
    \fi
}



\newenvironment{explanation}
{\textbf{解析:}\par}
{\par}



% 新增考察方面命令
\newcommand{\checklist}[1]{
    \textbf{考察:} #1 \par % 在题目下方显示考察方面
    \addcontentsline{toc}{subsection}{题号\thequestion 考察: #1} % 将考察内容添加到目录
    \vspace{2em} % 添加1em的垂直空间
}

\graphicspath{{images/}} %统一


\begin{document}
 %\fontsize{22}{31}\selectfont
 \tableofcontents % 添加目录
\newpage
%\pagestyle{empty}





\section{考点1:风险资本归属和风险调整后绩效衡量}
\setcounter{question}{0}
\begin{question}
    As part of its enterprise risk management enhancement initiative, Metropolitan Bank implemented a new economic capital-based decision framework that evaluates risk sources across various business divisions and organizational hierarchies. Which of the following statements about risk correlations is most accurate?\\
    作为其企业风险管理改进举措的一部分，大都会银行实施了一个新的经济资本为基础的决策框架，评估跨各种业务部门和组织层次的风险来源。以下关于风险相关性的陈述哪项最准确？
\end{question}

\begin{choices}
    \item Management decisions can influence the correlations between asset and liability risks.\\
    管理层决策可以影响资产和负债风险之间的相关性。
    \item Correlations between major risk categories (credit, market, operational) are well understood and easily measurable at the firm level.\\
    主要风险类别（信用、市场、操作）之间的相关性在公司层面得到很好的理解且容易衡量。
    \item Business unit correlations only matter for firm-wide economic capital allocation and are irrelevant for individual unit decisions.\\
    业务单元相关性仅对公司层面的经济资本分配重要，对单个单元决策无关。
    \item Including correlations in firm-wide risk assessment typically results in a total VaR that is greater than or equal to the sum of individual unit VaRs.\\
    在公司层面风险评估中包含相关性通常导致总VaR大于或等于单个单元VaR之和。
\end{choices}

\begin{correctanswer}
    A
\end{correctanswer}

\begin{explanation}
    \textbf{A选项正确。}
    \begin{itemize}
        \item 管理层决策可以影响资产和负债风险之间的相关性；例如，通过调整资产配置（如改变贷款组合、投资组合）、调整负债结构（如改变融资来源、期限匹配）、改变业务组合、实施对冲策略等方式，管理层可以影响资产和负债之间的相关性；这是风险管理的重要工具和策略。
    \end{itemize}

    \textbf{B选项错误。}
    \begin{itemize}
        \item 主要风险类别（信用、市场、操作）之间的相关性在公司层面并未得到很好理解且难以准确衡量；这些风险类别的数据来源、计量方法、时间范围不同，缺乏历史数据，且相关性在正常时期和压力时期可能不同；跨风险类别的相关性估计存在很大不确定性，是风险加总的主要挑战。
    \end{itemize}

    \textbf{C选项错误。}
    \begin{itemize}
        \item 业务单元相关性不仅对公司层面的经济资本配置重要，对单个单元决策也很重要；业务单元之间的相关性影响资源分配、风险限额设定、业务组合决策、绩效评估等；单个单元在做决策时需要考虑与其他单元的相关性，以优化整体风险回报。
    \end{itemize}

    \textbf{D选项错误。}
    \begin{itemize}
        \item 在公司层面风险评估中包含相关性通常导致总VaR小于（而非大于或等于）单个单元VaR之和；因为存在风险分散效应（diversification benefit），当不同单元或风险之间存在不完全正相关时，总风险会小于各单元风险的简单加总；只有当相关性为完全正相关（+1）时，总VaR才等于单个单元VaR之和。
    \end{itemize}
\end{explanation}

\checklist{经济资本框架中风险相关性的管理}


\begin{question}
    Which of the following statements about economic and regulatory capital are correct?
    \begin{itemize}
    \item I. Regulatory capital aims to maintain banking system stability by ensuring adequate capital levels
    \item II. Economic capital is designed to maintain solvency at a specific confidence level
    \item III. For any bank, economic capital must always be lower than regulatory capital
    \item IV. Economic capital allocation is a strategic process that impacts business unit and overall bank performance
    \end{itemize}
    以下关于经济资本和监管资本的陈述哪项正确？
    \begin{itemize}
        \item I. 监管资本旨在通过确保充足的资本水平来维持银行体系稳定性
        \item II. 经济资本旨在在特定置信水平下维持偿付能力
        \item III. 对于任何银行，经济资本必须始终低于监管资本
        \item IV. 经济资本配置是一个影响业务单元和整体银行绩效的战略过程
    \end{itemize}
\end{question}

\begin{choices}
    \item II and IV only
    \item I, II, III, and IV
    \item I, II, and IV only
    \item I and IV only
\end{choices}

\begin{correctanswer}
    C
\end{correctanswer}

\begin{explanation}
    \textbf{陈述I分析：正确}
    \begin{itemize}
        \item 监管资本旨在通过确保充足的资本水平来维持银行体系稳定；监管资本是监管机构要求银行持有的最低资本，目的是维护整个银行体系的稳定性、保护存款人利益、防止系统性风险，确保银行在压力时期仍能继续运营。
    \end{itemize}

    \textbf{陈述II分析：正确}
    \begin{itemize}
        \item 经济资本旨在维持特定置信水平下的偿付能力；经济资本是银行内部基于自身风险评估确定的资本，用于在特定置信水平（如99.9\%）下覆盖风险损失，确保银行在给定概率下保持偿付能力，是内部风险管理工具。
    \end{itemize}

    \textbf{陈述III分析：错误}
    \begin{itemize}
        \item 经济资本不一定始终低于监管资本；两者没有必然的大小关系；经济资本可能高于、等于或低于监管资本，取决于银行的风险状况、风险管理方法、监管要求、业务模式等因素；一些银行的经济资本可能高于监管资本（如果内部风险评估更保守），而另一些银行可能相反。
    \end{itemize}

    \textbf{陈述IV分析：正确}
    \begin{itemize}
        \item 经济资本分配是影响业务单元和整体银行绩效的战略过程；经济资本分配用于评估业务单元的风险调整后绩效（RAROC）、指导资源配置、设定风险限额、支持战略决策，直接影响业务单元和整体银行的风险调整后收益和绩效。
    \end{itemize}
\end{explanation}

\checklist{经济资本与监管资本的区别与联系}


\begin{question}
    Your bank calculates a one-day 95\% VaR for market risk, a one-year 99\% VaR for operational risk, and a one-year 99\% VaR for credit risk. The measures are \$80 million, \$400 million, and \$800 million, respectively. Operational risk is defined to include all risks that are not market risks and credit risks, and these three categories are mutually uncorrelated. The market risk VaR assumes normally distributed returns, and the bank expects to maintain its market risk VaR at that level for the whole year. Your supervisor wants your best estimate of a firm-wide VaR at the 1\% level. Among the following choices, your best estimate is:
    \\你的银行计算了一天的95\%市场风险VaR、一年的99\%操作风险VaR和一年的99\%信用风险VaR，分别为\$80 million、\$400 million和\$800 million。操作风险定义为不包括市场风险和信用风险的所有风险，这三个类别是相互独立的。市场风险VaR假设正态分布回报，银行期望全年保持该市场风险VaR水平。你的主管希望你给出公司层面的VaR在1\%置信水平下的最佳估计。在以下选项中，你的最佳估计是：
\end{question}

\begin{choices}
    \item \$1.36 billion
    \item \$1.55 billion
    \item \$2.00 billion
    \item It is impossible to aggregate risks with different distributions having only this information\\
    由于只有这些信息，无法将不同分布的风险加总
\end{choices}

\begin{correctanswer}
    C
\end{correctanswer}

\begin{explanation}
    \textbf{步骤1：将市场风险VaR转换为1年99\%水平}
    \begin{itemize}
        \item 市场风险：1天95\% VaR = \$80 million，需要转换为1年99\% VaR
        \item 置信水平转换：95\% → 99\%，乘数 = 2.326/1.645 ≈ 1.414（正态分布z值）
        \item 时间转换：1天 → 1年，乘数 = √252 ≈ 15.875（假设252个交易日，正态分布时间平方根规则）
        \item 调整后市场风险VaR = \$80 × 1.414 × 15.875 ≈ \$80 × 22.45 ≈ \$1,796 million
    \end{itemize}

    \textbf{步骤2：计算全公司VaR（三个风险类别不相关）}
    \begin{itemize}
        \item 由于三个风险类别相互不相关，使用平方根公式加总：
       \begin{align*}
       &\text{总VaR} = \sqrt{\text{市场风险VaR}^2 + \text{操作风险VaR}^2 + \text{信用风险VaR}^2} \\
       &= \sqrt{1{,}796^2 + 400^2 + 800^2} \\
       &= \sqrt{3{,}224{,}816 + 160{,}000 + 640{,}000} \\
       &= \sqrt{4{,}024{,}816} \approx \$2{,}006 \text{ million} \approx \$2{,}000 \text{ million}
       \end{align*}
    \end{itemize}
\end{explanation}

\checklist{不同风险类型VaR的加总方法}


\begin{question}
    The Chairman of a commercial bank recognizes that inadequate cash flow planning is one of the most critical mistakes a financial institution can make and insists that a robust liquidity-at-risk (LaR) measurement system be a core component of the bank's risk management framework. Which of the following statements about LaR is most accurate?\\
    一家商业银行的董事长认识到，资金流动性规划不足是金融机构可能犯的最严重错误之一，并坚持认为一个强大的流动性风险（LaR）测量系统应成为银行风险管理框架的核心组成部分。以下关于LaR的陈述哪项最准确？
\end{question}

\begin{choices}
    \item Hedging to reduce basis risk will lower LaR\\
    对冲以减少基差风险会降低LaR
    \item Futures hedging and long option positions have identical effects on LaR\\
    期货对冲和期权多头对LaR有相同的影响
    \item For a hedged portfolio, LaR can be substantially different from VaR\\
    对于对冲组合，LaR可能与VaR显著不同
    \item A bank's LaR typically decreases as its credit rating deteriorates\\
    银行信用评级下降时，LaR通常减少
\end{choices}

\begin{correctanswer}
    C
\end{correctanswer}

\begin{explanation}
    \textbf{A选项错误。}
    \begin{itemize}
        \item 对冲以减少基差风险不一定会降低LaR；对冲基差风险可能涉及期货或衍生品，这些工具会产生保证金要求，增加流动性需求；虽然对冲可能降低市场风险（VaR），但可能增加流动性风险（LaR），因为保证金要求会产生现金流流出。
    \end{itemize}

    \textbf{B选项错误。}
    \begin{itemize}
        \item 期货对冲和期权多头对LaR的影响不同；期货对冲需要初始保证金和维持保证金，会产生持续的现金流流出，增加流动性风险；期权多头通常只需要支付权利金，没有后续保证金要求，对LaR的影响较小；两者对LaR的影响显著不同。
    \end{itemize}

    \textbf{C选项正确。}
    \begin{itemize}
        \item 对于对冲组合，LaR可能与VaR显著不同；VaR衡量市场价值损失风险，而LaR衡量流动性需求/现金流风险；对冲组合可能降低VaR（因为市场风险被对冲），但可能增加LaR（因为期货对冲需要保证金，产生流动性需求）；例如，期货对冲可能降低VaR但增加保证金要求，欧式期权可能降低VaR且不影响LaR（如果已支付权利金），两者可能存在显著差异。
    \end{itemize}

    \textbf{D选项错误。}
    \begin{itemize}
        \item 银行信用评级下降时，LaR通常增加而非减少；信用评级下降会导致融资成本上升、融资渠道受限、市场信心下降、存款流失风险增加，这些都会增加流动性压力和流动性风险，因此LaR通常会增加。
    \end{itemize}
\end{explanation}

\checklist{流动性风险与市场风险的关系}



\begin{question}
    You are a portfolio manager at a leading investment fund analyzing a 2,000 share position in an undervalued but illiquid stock XYZ, which has a current stock price of USD 85 (expressed as the midpoint of the current bid-ask spread.) Daily return for XYZ has an estimated volatility of 1.15\%. The average bid-ask spread is USD 0.18. Assuming returns of XYZ are normally distributed, what is the estimated liquidity-adjusted daily 95\% VaR, using the constant spread approach?\\
    你是一家领先投资基金的组合经理，分析一个2,000股头寸的XYZ股票，该股票估值过低但流动性不足，当前股票价格为USD 85（表示为当前买卖价差的中间价）。XYZ的日回报估计波动率为1.15\%，平均买卖价差为USD 0.18。假设XYZ的回报呈正态分布，使用固定价差方法估计流动性调整后的每日95\% VaR是多少？
\end{question}

\begin{choices}
    \item USD 2,450
    \item USD 2,580
    \item USD 2,710
    \item USD 2,840
\end{choices}

\begin{correctanswer}
    C
\end{correctanswer}

\begin{explanation}
    \textbf{步骤1：计算头寸价值}
    \begin{itemize}
        \item 头寸价值 = 2,000股 × \$85/股 = \$170,000
    \end{itemize}

    \textbf{步骤2：计算市场风险VaR（不考虑流动性）}
    \begin{itemize}
        \item 市场风险VaR公式：$VaR = \text{头寸价值} \times z_{\alpha} \times \sigma$
        \item 其中：$z_{0.95} = 1.645$（95\%置信水平），$\sigma = 1.15\% = 0.0115$（日波动率）
        \item 市场风险VaR = \$170,000 × (1.645 × 0.0115) = \$170,000 × 0.0189175 = \$2,580
    \end{itemize}

    \textbf{步骤3：计算流动性成本（固定价差方法）}
    \begin{itemize}
        \item 流动性成本公式：$\text{流动性成本} = \text{头寸价值} \times \left(\frac{1}{2} \times \frac{\text{买卖价差}}{\text{价格}}\right)$
        \item 流动性成本 = \$170,000 × (0.5 × 0.18/85) = \$170,000 × 0.001059 = \$130
        \item 注：0.5表示买卖价差的一半（买入或卖出时承担的一半价差成本）
        \item 流动性调整VaR = 市场风险VaR + 流动性成本 = \$2,580 + \$130 = \$2,710

    \end{itemize}
   
\end{explanation}

\checklist{流动性调整VaR的计算方法}


\begin{question}
    Global Capital is an investment management firm with USD 30 billion under management. It holds 25\% of the shares in a company. Global Capital's risk manager is concerned that, if the entire position needs to be liquidated, its size would impact the market price. The estimated price elasticity of demand is -0.4. What is the increase in Global Capital's Value-at-Risk estimate for this position if a liquidity adjustment is made?\\
    Global Capital是一家管理USD 30 billion的投资管理公司，持有某公司25\%的股份。Global Capital的风险经理担心，如果整个头寸需要清算，其规模会影响市场价格。估计的需求价格弹性为-0.4。如果进行流动性调整，Global Capital对该头寸的Value-at-Risk估计增加多少？
\end{question}

\begin{choices}
    \item 5%
    \item 10%
    \item 15%
    \item 20%
\end{choices}

\begin{correctanswer}
    B
\end{correctanswer}

\begin{explanation}
    \textbf{步骤1：理解题目条件}
    \begin{itemize}
        \item Global Capital持有该公司25\%的股份，即交易规模比例 = 0.25
        \item 价格弹性需求 = -0.4（负号表示价格上升时需求下降）
        \item 需要计算流动性调整后VaR相对于原始VaR的增加比例
    \end{itemize}

    \textbf{步骤2：计算价格变化率}
    \begin{itemize}
        \item 价格弹性公式：$\text{价格变化率} = \text{弹性} \times \text{交易规模比例}$
        \item $\Delta P/P = -0.4 \times 0.25 = -0.1 = -10\%$
        \item 这意味着如果全部头寸需要清算，价格将下降10\%
    \end{itemize}

    \textbf{步骤3：计算流动性调整VaR}
    \begin{itemize}
        \item 流动性调整VaR与原始VaR的比例公式：$LVaR/VaR = 1 - (\Delta P/P)$
        \item $LVaR/VaR = 1 - (-0.1) = 1 + 0.1 = 1.1$
        \item 这表示流动性调整VaR是原始VaR的1.1倍
    \end{itemize}

    \textbf{步骤4：计算VaR增加比例}
    \begin{itemize}
        \item VaR增加比例 = $(LVaR - VaR)/VaR = 1.1 - 1 = 0.1 = 10\%$
        \item 因此，流动性调整使VaR增加了10\%
    \end{itemize}
\end{explanation}

\checklist{流动性调整对VaR的影响计算}


\begin{question}
    In calculating its risk-adjusted return on capital, your bank uses a capital charge of 3.00\% for revolving credit facilities with a loan equivalent factor of 0.40 assigned to the undrawn portion. Recently, you have become concerned that the protective covenants in these loans are insufficient and may not prevent customers from drawing on the facilities during stress periods. As such, you have recommended increasing the loan equivalent factor to 0.80. This recommendation has met with resistance from the loan origination team, and senior management has asked you to quantify the impact of your recommendation. For a typical facility that has an original principal of USD 1.2 billion and is 25\% drawn, how much additional economic capital would have to be allocated if you increase the loan equivalent factor from 0.40 to 0.80?\\
    在计算风险调整后资本回报率时，你的银行对未提取部分的贷款等效因子为0.40的循环信贷设施使用3.00\%的资本费用。最近，你担心这些贷款的保护性条款不足，可能无法在压力期间阻止客户提取资金。因此，你建议增加贷款等效因子至0.80。这个建议遇到了贷款发起团队的抵制，高级管理层要求你量化你的建议的影响。对于一个原始本金为USD 1.2 billion且已提取25\%的典型设施，如果你将贷款等效因子从0.40增加到0.80，需要分配多少额外的经济资本？
\end{question}

\begin{choices}
    \item USD 4.32 million
    \item USD 7.20 million
    \item USD 10.80 million
    \item USD 15.12 million
\end{choices}

\begin{correctanswer}
    B
\end{correctanswer}

\begin{explanation}
   

    \textbf{B选项正确。}
    \begin{itemize}
        \item 初始资本要求 = [(12亿 × 0.25) + (12亿 × 0.75 × 0.40)] × 3\% = 16.2百万
        \item 调整后资本要求 = [(12亿 × 0.25) + (12亿 × 0.75 × 0.80)] × 3\% = 23.4百万
        \item 额外资本要求 = 23.4 - 16.2 = 7.2百万
    \end{itemize}

\end{explanation}

\checklist{循环信贷额度经济资本计算}


\begin{question}
    Michael Chen, an analyst at Global Advisors, is analyzing the economic effects of buying stock with borrowed funds for a high net worth individual client. Assume that the client has \$300 cash invested (i.e., no borrowed funds) and then uses the cash to purchase stock. The client then decides to use 60\% borrowed funds to purchase stock on margin. After the margin transaction, the total assets on the full economic balance sheet and the leverage ratio are closest to:\\
    Michael Chen，Global Advisors的分析师，正在分析高净值个人客户使用借入资金购买股票的经济影响。假设客户有\$300现金投资（即没有借入资金），然后使用现金购买股票。客户随后决定使用60\%借入资金以保证金方式购买股票。在保证金交易后，全经济资产负债表上的总资产和杠杆率最接近：
\end{question}

\begin{choices}
    \item Total assets: \$300; Leverage ratio: 1.6\\
    总资产: \$300; 杠杆率: 1.6
    \item Total assets: \$450; Leverage ratio: 1.6\\
    总资产: \$450; 杠杆率: 1.6
    \item Total assets: \$450; Leverage ratio: 2.5\\
    总资产: \$450; 杠杆率: 2.5
    \item Total assets: \$600; Leverage ratio: 2.5\\
    总资产: \$600; 杠杆率: 2.5
\end{choices}

\begin{correctanswer}
    B
\end{correctanswer}

\begin{explanation}
  

    \textbf{B选项正确。}
    \begin{itemize}
        \item 初始状态：现金\$300，无负债
        \item 保证金交易：借入\$180（60\%）
        \item 总资产 = \$300 + \$180 = \$450
        \item 杠杆率 = 总资产/权益 = \$450/\$300 = 1.6
    \end{itemize}

   
\end{explanation}

\checklist{保证金交易的杠杆效应计算}


\begin{question}
    A trader observes a quote for Stock ABC, and the midpoint of its current best bid and best ask price is CAD 42. ABC has an estimated daily return volatility of 0.30\% and average bid-ask spread of CAD 0.12\%. Assuming the returns of ABC are normally distributed, what is closest to the estimated liquidity-adjusted, 1-day 95\% VaR, using the constant spread approach on a 15,000 share position?\\
    一个交易员观察到Stock ABC的报价，其当前最佳买卖价中间价为CAD 42。ABC的日回报估计波动率为0.30\%，平均买卖价差为CAD 0.12\%。假设ABC的回报呈正态分布，使用固定价差方法估计流动性调整后的1天95\% Va
\end{question}

\begin{choices}
    \item CAD 2,500
    \item CAD 3,200
    \item CAD 3,900
    \item CAD 4,600
\end{choices}

\begin{correctanswer}
    B
\end{correctanswer}

\begin{explanation}
   

    \textbf{B选项正确。}
    这个说法正确：
    \begin{itemize}
        \item 市场风险VaR = 42 × 15,000 × (1.645 × 0.003) = 2,700
        \item 流动性成本 = 630,000 × (0.5 × 0.12/42) = 500
        \item 流动性调整VaR = 2,700 + 500 = 3,200
    \end{itemize}

   
\end{explanation}

\checklist{流动性调整VaR的计算方法}


\begin{question}
    A firm has determined that the risk-adjusted return on capital (RAROC) for a particular project is 16\%. To evaluate whether the firm should accept the project, an analyst determines that the firm's beta is 1.4, the expected market return is 14\%, and the risk-free interest rate is 6\%. If the analyst uses the adjusted RAROC (ARAROC) methodology to make an accept/reject decision, should the project be accepted?
    \\一个公司确定了一个特定项目的风险调整后资本回报率（RAROC）为16\%。为了评估公司是否应该接受该项目，分析师确定公司的贝塔系数为1.4，预期市场回报率为14\%，无风险利率为6\%。如果分析师使用调整后的RAROC（ARAROC）方法做出接受/拒绝决策，该项目是否应该被接受？
\end{question}

\begin{choices}
    \item No, because the computed ARAROC is approximately 1\% less than the market risk premium\\
    不接受，因为计算的ARAROC大约比市场风险溢价低1\%
    \item No, because the RAROC is 1.25\% less than the return predicted by the CAPM\\
    不接受，因为RAROC比CAPM预测的回报低1.25\%
    \item Yes, because the computed ARAROC is approximately 1\% more than the market risk premium\\
    接受，因为计算的ARAROC大约比市场风险溢价高1\%
    \item Yes, because the ARAROC is approximately 4\% more than the return predicted by the CAPM\\
    接受，因为计算的ARAROC大约比市场风险溢价高4\%
\end{choices}

\begin{correctanswer}
    A
\end{correctanswer}

\begin{explanation}
    \textbf{步骤1：计算ARAROC}
    \begin{itemize}
        \item ARAROC公式：$ARAROC = \frac{RAROC - R_f}{\beta}$
        \item 其中：RAROC = 16\%，$R_f$ = 6\%，$\beta$ = 1.4
        \item ARAROC = (0.16 - 0.06) / 1.4 = 0.10 / 1.4 = 7.14\%
    \end{itemize}

    \textbf{步骤2：计算市场风险溢价}
    \begin{itemize}
        \item 市场风险溢价 = 预期市场回报 - 无风险利率
        \item 市场风险溢价 = 14\% - 6\% = 8\%
    \end{itemize}

    \textbf{步骤3：比较ARAROC与市场风险溢价}
    \begin{itemize}
        \item ARAROC (7.14\%) < 市场风险溢价 (8\%)
        \item 差异 = 8\% - 7.14\% = 0.86\% $\approx$ 1\%
        \item ARAROC比市场风险溢价低约1\%
    \end{itemize}

    \textbf{步骤4：决策规则}
    \begin{itemize}
        \item ARAROC决策规则：ARAROC应至少等于市场风险溢价才能接受项目
        \item 因为ARAROC (7.14\%) < 市场风险溢价 (8\%)，项目提供的风险调整后回报不足以补偿市场风险，因此应拒绝该项目
    \end{itemize}

\end{explanation}

\checklist{ARAROC与市场风险溢价的比较}


\begin{question}
    A bank issues a \$250,000,000 loan with the following characteristics:
    \begin{itemize} 
        \item Loan pays a fixed annual interest rate of 9.0\%
        \item The interest expense associated with the loan is 6.5\%
        \item The operating cost to the bank's commercial lending division is \$2,000,000
        \item Economic capital required to support the loan is \$20 million, which is invested in T-bills paying a rate of 3.0\%
        \item The expected loss for the loan is 20 basis points per year
    \end{itemize}
    What is the risk-adjusted return on capital (RAROC) for this loan?
    \\一家银行发行了一个\$250,000,000贷款，具有以下特征：
    \begin{itemize}
        \item 贷款支付固定年利率9.0\%
        \item 与贷款相关的利息费用为6.5\%
        \item 银行商业贷款部门的运营成本为\$2,000,000
        \item 支持贷款的经济资本为\$20 million，投资于T-bills支付3.0\%的利率
        \item 贷款的预期损失为每年20个基点
    \end{itemize}
    该贷款的风险调整后资本回报率（RAROC）是多少？
\end{question}

\begin{choices}
    \item 3.50\%
    \item 7.25\%
    \item 19.50\%
    \item 23.50\%
\end{choices}

\begin{correctanswer}
    D
\end{correctanswer}

\begin{explanation}
  
    \textbf{D选项正确。}
    \begin{itemize}
        \item 收入 = 250,000,000 × 9\% = 22,500,000
        \item 预期损失 = 250,000,000 × 0.2\% = 500,000
        \item 费用 = 250,000,000 × 6.5\% + 2,000,000 = 18,250,000
        \item 经济资本收益 = 20,000,000 × 3\% = 600,000
        \item RAROC = (22,500,000 - 500,000 - 18,250,000 + 600,000)/20,000,000 = 23.5\%
    \end{itemize}
\end{explanation}

\checklist{RAROC的计算方法}


\begin{question}
    Suppose a bank loan has the following characteristics:
    \begin{itemize}
        \item Gross revenues are \$1.5 million
        \item Interest expense is \$0.4 million
        \item The return on the \$6 million economic capital invested in T-bills is \$150,000
        \item The firm beta is 1.6
        \item The unexpected loss for the loan is estimated at \$600,000
        \item The adjusted RAROC is equal to 9\%
    \end{itemize}
    What is the worst-case loss for this loan?
    \\假设一个银行贷款具有以下特征：
    \begin{itemize}
        \item 总收入为\$1.5 million
        \item 利息费用为\$0.4 million
        \item 投资于T-bills的\$6 million经济资本回报为\$150,000
        \item 公司贝塔系数为1.6
        \item 贷款的意外损失估计为\$600,000
        \item 调整后的RAROC为9\%
    \end{itemize}
    该贷款的最坏情况损失是多少？
\end{question}

\begin{choices}
    \item \$1.8 million
    \item \$2.5 million
    \item \$1.1 million
    \item \$0.85 million
\end{choices}

\begin{correctanswer}
    C
\end{correctanswer}

\begin{explanation}
   
    \textbf{C选项正确。}
    \begin{itemize}
        \item 无风险利率 = 150,000/6,000,000 = 2.5\%
        \item RAROC = 1.6 × 0.09 + 0.025 = 16.9\%
        \item 预期损失 = -(0.169 × 6M - 1.5M + 0.4M - 0.15M) = 500,000
        \item 最坏情况损失 = 500,000 + 600,000 = 1.1M
    \end{itemize}

\end{explanation}

\checklist{最坏情况损失的计算方法}

\begin{question}
    In its efforts to enhance its enterprise risk management function, Metropolitan Bank implemented a new economic capital-based decision framework that evaluates risk sources across various business divisions and organizational hierarchies. Which of the following statements about risk correlations is most accurate?
    \\为了增强其企业风险管理功能，大都会银行实施了一个新的经济资本为基础的决策框架，评估跨各种业务部门和组织层次的风险来源。以下关于风险相关性的陈述哪项最准确？
\end{question}

\begin{choices}
    \item Management decisions can influence the correlations between asset and liability risks\\
    管理层决策可以影响资产和负债风险之间的相关性
    \item Correlations between major risk categories (credit, market, operational) are well understood and easily measurable at the firm level\\
    主要风险类别（信用、市场、操作）之间的相关性在公司层面得到很好的理解且容易衡量
    \item Business unit correlations only matter for firm-wide economic capital allocation and are irrelevant for individual unit decisions\\
    业务单元相关性仅对公司层面的经济资本分配重要，对单个单元决策无关
    \item Including correlations in firm-wide risk assessment typically results in a total VaR that is greater than or equal to the sum of individual unit VaRs\\
    在公司层面风险评估中包含相关性通常导致总VaR大于或等于单个单元VaR之和
\end{choices}

\begin{correctanswer}
    A
\end{correctanswer}

\begin{explanation}
    \textbf{A选项正确。}
    \begin{itemize}
        \item 管理层决策可以影响资产和负债风险之间的相关性；例如，通过调整资产配置（如改变贷款组合、投资组合）、调整负债结构（如改变融资来源、期限匹配）、改变业务组合、实施对冲策略等方式，管理层可以影响资产和负债之间的相关性；这是风险管理的重要工具和策略。
    \end{itemize}

    \textbf{B选项错误。}
    \begin{itemize}
        \item 主要风险类别（信用、市场、操作）之间的相关性在公司层面并未得到很好理解且难以准确衡量；这些风险类别的数据来源、计量方法、时间范围不同，缺乏历史数据，且相关性在正常时期和压力时期可能不同；跨风险类别的相关性估计存在很大不确定性，是风险加总的主要挑战。
    \end{itemize}

    \textbf{C选项错误。}
    \begin{itemize}
        \item 业务单元相关性不仅对公司层面的经济资本配置重要，对单个单元决策也很重要；业务单元之间的相关性影响资源分配、风险限额设定、业务组合决策、绩效评估等；单个单元在做决策时需要考虑与其他单元的相关性，以优化整体风险回报。
    \end{itemize}

    \textbf{D选项错误。}
    \begin{itemize}
        \item 在公司层面风险评估中包含相关性通常导致总VaR小于（而非大于或等于）单个单元VaR之和；因为存在风险分散效应（diversification benefit），当不同单元或风险之间存在不完全正相关时，总风险会小于各单元风险的简单加总；只有当相关性为完全正相关（+1）时，总VaR才等于单个单元VaR之和。
    \end{itemize}
\end{explanation}

\checklist{经济资本框架中风险相关性的管理}


\begin{question}
    A treasury analyst at a bank is using the RAROC approach to assess whether to extend a new CNY 1 billion loan. The loan would be fully funded by deposits and would have the following metrics:
    \begin{itemize}
        \item Required economic capital: CNY 120 million
        \item Interest rate on the loan: 5.5\%
        \item Operating expenses to support the loan: CNY 8 million
        \item Expected loss on the loan: CNY 15 million
        \item Unexpected loss on the loan: CNY 12 million
        \item Transfer price charge on deposits funding the loan: 1.2\%
        \item Effective tax rate: 25\%
    \end{itemize}
    The analyst assumes a return of 2.5\% on the bank's economic capital over the duration of the loan. What is the after-tax RAROC for this loan?
    \\一家银行的 treasury analyst正在使用RAROC方法评估是否延长一个新的人民币10亿元贷款。贷款将完全由存款支持，具有以下指标：
    \begin{itemize}
        \item 所需经济资本：人民币120亿元
        \item 贷款利率：5.5\%
        \item 支持贷款的运营成本：人民币8亿元
        \item 贷款预期损失：人民币15亿元
        \item 贷款意外损失：人民币12亿元
        \item 存款支持贷款的转移价格费用：1.2\%
        \item 有效税率：25\%
    \end{itemize}
    分析师假设贷款期间银行经济资本的回报率为2.5\%。该贷款的税后RAROC是多少？
\end{question}

\begin{choices}
    \item 5.5\%
    \item 6.2\%
    \item 10.8\%
    \item 12.5\%
\end{choices}

\begin{correctanswer}
    D
\end{correctanswer}

\begin{explanation}

    \textbf{D选项正确。}
    \begin{itemize}
        \item 收入 = 10亿 × 5.5\% = 5,500万
        \item 费用 = 10亿 × 1.2\% + 800万 = 2,000万
        \item 预期损失 = 1,500万
        \item 经济资本收益 = 1.2亿 × 2.5\% = 300万
        \item 税前收益 = 5,500万 - 2,000万 - 1,500万 + 300万 = 2,300万
        \item 税后RAROC = 2,300万 × (1-25\%)/1.2亿 = 12.5\%
    \end{itemize}
\end{explanation}

\checklist{税后RAROC的计算方法}





\newpage


\section{考点2:经济资本框架中的各种做法和问题}
\setcounter{question}{0}
\begin{question}
    Which of the following statements about the differences between market and operational value-at-risk at financial institutions are correct?
    \begin{itemize}
    \item I. The distribution of operational risk events must include sufficient mass in the extreme tail, making an assumption of a normal distribution invalid.
    \item II. The typical time horizon of market VaR calculations is 1 day, whereas the typical time horizon of operational VaR calculations is 1 year.
    \item III. Since prices are sufficiently available for liquid assets at all times, the market risk of liquid assets can be modeled using continuous distributions, but the nature of operational risk events requires using discrete distributions.
    \item IV. Market VaR requires a higher confidence level than operational VaR.
    \end{itemize}
    以下关于金融机构市场风险和操作风险VaR区别的陈述哪项正确？
    \begin{itemize}
        \item I. 操作风险事件分布必须在极端尾部有足够的质量，使得正态分布假设无效。
        \item II. 市场VaR计算的典型时间范围是1天，而操作VaR计算的典型时间范围是1年。
        \item III. 由于流动性资产的价格足够可得，市场风险可以建模为连续分布，但操作风险事件的性质需要使用离散分布。
        \item IV. 市场VaR需要比操作VaR更高的置信水平。
    \end{itemize}
\end{question}

\begin{choices}
    \item I, II, and III
    \item I, II and IV
    \item I, II, III and IV
    \item III and IV
\end{choices}

\begin{correctanswer}
    A
\end{correctanswer}

\begin{explanation}
    \textbf{陈述I分析：正确}
    \begin{itemize}
        \item 操作风险事件分布必须在极端尾部有足够的质量，使得正态分布假设无效；操作风险具有厚尾特征，极端损失事件（低频高损）的概率显著高于正态分布预测，因此不能用正态分布建模，需要使用重尾分布（如对数正态、帕累托、广义帕累托等）。
    \end{itemize}

    \textbf{陈述II分析：正确}
    \begin{itemize}
        \item 市场VaR计算的典型时间范围是1天，而操作VaR计算的典型时间范围是1年；市场风险变化频繁，需要短期度量（1天），而操作风险事件频率较低，需要长期度量（1年）以捕获足够的事件；这是两种VaR的主要区别之一。
    \end{itemize}

    \textbf{陈述III分析：正确}
    \begin{itemize}
        \item 由于流动性资产的价格足够可得，市场风险可以建模为连续分布，但操作风险事件的性质需要使用离散分布；市场风险资产价格连续变化，可用连续分布；操作风险事件的发生是离散的（事件要么发生要么不发生），频率建模使用离散分布（如泊松、负二项），严重度可能使用连续分布，但整体事件性质是离散的。
    \end{itemize}

    \textbf{陈述IV分析：错误}
    \begin{itemize}
        \item 市场VaR不需要比操作VaR更高的置信水平；实际上，操作VaR通常使用更高的置信水平（如99.9\%用于监管资本），而市场VaR通常使用较低的置信水平（如95\%或99\%）；置信水平是用户设定的参数，没有必然的高低关系。
    \end{itemize}
\end{explanation}

\checklist{市场风险与操作风险VaR的区别}



\begin{question}
    Your Chief Risk Officer asks you to prepare a list of early warning indicators for liquidity problems for your bank. Which of the following are early warning indicators of a potential liquidity problem?
    \begin{itemize}
    \item I. Rapid asset growth, especially when funded with potentially volatile liabilities.
    \item II. Growing concentrations in assets or liabilities.
    \item III. An increase of the weighted average maturity of liabilities.
    \item IV. Reduction in the frequency of positions approaching or breaching internal or regulatory limits.
    \item V. Narrowing debt or credit default swap spreads.
    \item VI. Counterparties that request additional collateral for credit exposures.
    \item VII. Increasing redemptions of CDs before maturity.
    \end{itemize}
    你的首席风险官要求你准备一份流动性问题早期预警指标列表。以下哪些是潜在流动性问题的早期预警指标？
    \begin{itemize}
        \item I. 快速资产增长，特别是当用潜在不稳定负债融资时。
        \item II. 资产或负债集中度增加。
        \item III. 负债加权平均期限增加。
        \item IV. 接近或突破内部或监管限额的频率降低。
        \item V. 债务或信用违约互换利差收窄。
        \item VI. 交易对手要求增加信用敞口的抵押品。
        \item VII. 提前赎回的CD数量增加。
    \end{itemize}
\end{question}

\begin{choices}
    \item I, II, VI, and VII
    \item I, III, V, and VI
    \item II, IV, V, and VII
    \item I, V, VI, and VII
\end{choices}

\begin{correctanswer}
    A
\end{correctanswer}

\begin{explanation}
    \textbf{陈述I分析：是}
    \begin{itemize}
        \item 快速资产增长，特别是用潜在不稳定负债融资，是流动性问题的早期预警指标；资产快速增长如果依赖短期或波动性负债（如批发融资、回购协议），可能导致流动性错配，在压力时期这些资金来源可能迅速撤离，引发流动性危机。
    \end{itemize}

    \textbf{陈述II分析：是}
    \begin{itemize}
        \item 资产或负债集中度增加是流动性问题的早期预警指标；集中度增加意味着对特定资金来源、资产类型或交易对手的依赖增加，流动性风险增大；如果主要资金来源或资产类型面临压力，银行可能面临重大流动性问题。
    \end{itemize}

    \textbf{陈述III分析：不是}
    \begin{itemize}
        \item 负债加权平均期限增加通常不是流动性问题指标，反而是流动性改善的信号；长期负债提供更稳定的资金来源，降低短期融资需求和流动性风险；期限延长通常有助于降低流动性风险。
    \end{itemize}

    \textbf{陈述IV分析：不是}
    \begin{itemize}
        \item 接近或突破内部或监管限额的频率降低不是流动性问题指标，反而是风险改善的信号；这表明银行更好地管理风险，保持在限额内，而非流动性问题的早期预警。
    \end{itemize}

    \textbf{陈述V分析：不是}
    \begin{itemize}
        \item 债务或信用违约互换利差收窄通常不是流动性问题指标；利差收窄通常表示信用状况改善、流动性改善或市场信心增强，而非流动性问题的早期预警；流动性问题通常伴随着利差扩大。
    \end{itemize}

    \textbf{陈述VI分析：是}
    \begin{itemize}
        \item 交易对手要求增加信用敞口的抵押品是流动性问题的早期预警指标；这表明市场对银行的信用和流动性担忧增加，交易对手要求更多保护，可能预示流动性压力；这也会直接消耗银行的流动性资产。
    \end{itemize}
\end{explanation}

\checklist{流动性风险的早期预警指标}


\begin{question}
    Metropolitan Bank approaches economic capital and risk aggregation by first estimating the stand-alone economic capital for individual risk factors. In a second step, the bank aggregates risks based on the relative amounts of economic capital allocated to these risks, taking into account the correlations between risk factors. Which of the following variables is not a primary driver of the diversification benefit that accrues from aggregation?
    \\大都会银行首先估计单个风险因素的经济资本，然后在第二步中，根据分配给这些风险的相对经济资本量，考虑风险因素之间的相关性，对风险进行聚合。以下哪个变量不是分散化收益的主要驱动因素？
\end{question}

\begin{choices}
    \item The number of risk positions\\
    风险头寸数量
    \item The size of the portfolio\\
    组合规模
    \item The concentration of those risk positions, or their relative weights in a portfolio\\
    风险头寸集中度或相对权重
    \item The correlation between the positions\\
    头寸间相关性
\end{choices}

\begin{correctanswer}
    B
\end{correctanswer}

\begin{explanation}
    \textbf{A选项是主要驱动因素。}
    \begin{itemize}
        \item 风险头寸数量是分散化收益的主要驱动因素；在相关性相同的情况下，头寸数量越多，分散化效果越好；这是因为风险分散的程度取决于头寸的多样性和数量，更多头寸提供更多分散化机会（前提是相关性不完全正相关）。
    \end{itemize}

    \textbf{B选项不是主要驱动因素（正确答案）。}
    \begin{itemize}
        \item 组合规模不是分散化收益的主要驱动因素；分散化收益通常以比例或百分比表示，与绝对规模无关；如果组合规模翻倍，总风险也大约翻倍，但分散化收益的比例保持相对稳定；分散化收益主要取决于头寸的相对权重、相关性和数量结构，而非绝对规模。
    \end{itemize}

    \textbf{C选项是主要驱动因素。}
    \begin{itemize}
        \item 风险头寸集中度或相对权重是分散化收益的主要驱动因素；集中度越低（即头寸权重越均衡），分散化效果越好；如果组合高度集中在少数头寸，分散化收益有限；均衡的权重分布有助于最大化分散化收益。
    \end{itemize}

    \textbf{D选项是主要驱动因素。}
    \begin{itemize}
        \item 头寸间相关性是分散化收益的主要驱动因素；相关性越低，分散化效果越好；当相关性接近1时，分散化收益最小；当相关性接近0或为负时，分散化收益最大；相关性是决定分散化收益的关键因素。
    \end{itemize}
\end{explanation}

\checklist{风险分散化收益的驱动因素}


\begin{question}
    Consider a bank that wants to have an amount of capital so that it can absorb unexpected losses corresponding to a firm-wide VaR at the 1\% level. It measures firm-wide VaR by adding up the VaRs for market risk, operational risk, and credit risk. There is a risk that the bank has too little capital because\\
    考虑一家银行希望拥有足够的资本，以吸收与公司层面VaR（1\%置信水平）相对应的意外损失。它通过将市场风险、操作风险和信用风险的VaR相加来衡量公司层面VaR。存在银行资本不足的风险，因为
\end{question}

\begin{choices}
    \item It does not take into account the correlations among risks\\
    它没有考虑风险之间的相关性
    \item It ignores risks that are not market, operational, or credit risks\\
    它忽略了不是市场、操作或信用风险的其他风险
    \item It mistakenly uses VaR to measure operational risk because operational risks that matter are rare events\\
    它错误地使用VaR衡量操作风险，因为重要的操作风险是罕见事件
    \item It is meaningless to add VaRs\\
    加总VaR没有意义
\end{choices}

\begin{correctanswer}
    B
\end{correctanswer}

\begin{explanation}
    \textbf{A选项错误。}
    \begin{itemize}
        \item 不考虑风险之间的相关性不会导致资本不足，反而可能导致资本要求高估；如果忽略相关性（即假设完全正相关），简单相加VaR会高估总风险，导致资本要求过高；题目问的是"资本太少"的原因，而非太多。
    \end{itemize}

    \textbf{B选项正确。}
    \begin{itemize}
        \item 忽略非市场、操作、信用风险的其他风险会导致资本不足；如果仅考虑三大风险类别，可能忽略了其他重要风险类型（如流动性风险、集中度风险、战略风险、声誉风险、模型风险等），这些风险未被纳入全公司VaR计算，导致总风险低估，从而资本要求不足；这是导致资本不足的主要原因。
    \end{itemize}

    \textbf{C选项错误。}
    \begin{itemize}
        \item 使用VaR衡量操作风险不是导致资本不足的原因；虽然操作风险涉及罕见事件，但如果VaR置信水平设置得当（如99\%用于监管资本），应该能覆盖这些罕见事件；如果VaR计算正确，即使操作风险是罕见事件，也应该在VaR中反映；这不是导致资本不足的原因。
    \end{itemize}

    \textbf{D选项错误。}
    \begin{itemize}
        \item 加总VaR并非没有意义，虽然加总可能不够精确，但这通常会导致资本要求高估而非低估（因为没有考虑分散化效应）；如果简单相加VaR，假设完全正相关，会高估总风险，导致资本要求过高；这不是导致资本不足的原因。
    \end{itemize}
\end{explanation}

\checklist{银行资本充足性的风险覆盖范围}


\begin{question}
    Based on "Supervisory Guidance for Assessing Banks' Financial Instrument Fair Value Practices" issued by the Basel Committee, which of the following factors should be considered in determining whether the sources of fair values are reliable and relevant?
\begin{itemize}
    \item I. Frequency and availability of prices/quotes
    \item II. Maturity of the market
    \item III. Agreement of values with those generated by internal models
    \item IV. Number of independent sources that produce the prices/quotes
\end{itemize}
根据巴塞尔委员会发布的“评估银行金融工具公允价值实践的监管指导”，在确定公允价值来源是否可靠和相关时，应考虑以下哪些因素？
    \begin{itemize}
        \item I. 价格/报价的频率和可得性
        \item II. 市场的成熟度
        \item III. 与内部模型生成的价值的一致性
        \item IV. 产生价格/报价的独立来源数量
    \end{itemize}
\end{question}

\begin{choices}
    \item I and II only
    \item III and IV only
    \item I, II and III only
    \item I, II and IV only
\end{choices}

\begin{correctanswer}
    D
\end{correctanswer}

\begin{explanation}
    \textbf{陈述I分析：应该考虑}
    \begin{itemize}
        \item 价格/报价的频率和可得性是确定公允价值来源可靠性和相关性的重要因素；高频率和良好的可得性表明市场活跃、流动性好，价格能够及时反映市场信息，提高公允价值的可靠性和相关性。
    \end{itemize}

    \textbf{陈述II分析：应该考虑}
    \begin{itemize}
        \item 市场成熟度是确定公允价值来源可靠性和相关性的重要因素；成熟的市场通常具有更好的流动性、更透明的价格发现机制和更稳定的交易环境，这些特征提高了价格数据的可靠性和相关性。
    \end{itemize}

    \textbf{陈述III分析：不应该考虑}
    \begin{itemize}
        \item 与内部模型生成的价值的一致性不是确定公允价值来源可靠性和相关性的必要因素；公允价值应以市场价格为准，而非内部模型；如果市场价格与内部模型不一致，可能说明市场价格本身是可靠的，而用内部模型来验证市场价格的做法不符合公允价值的基本原则；巴塞尔指导强调市场数据本身的特征，而非与内部模型的一致性。
    \end{itemize}

    \textbf{陈述IV分析：应该考虑}
    \begin{itemize}
        \item 产生价格/报价的独立来源数量是确定公允价值来源可靠性和相关性的重要因素；多个独立来源可以提高价格数据的可靠性，减少单一来源的偏差或错误风险，多个来源的一致性验证可以提高数据的可信度。
    \end{itemize}
\end{explanation}

\checklist{公允价值可靠性的评估因素}


\begin{question}
    Alpha Capital, LLC (Alpha) uses monthly model vetting to mitigate potential model risk. Alpha's managers recently accepted the use of a model for valuing short-term options on 25-year corporate bonds, but rejected the same model to value short-term options on five-year government bonds. The managers also frequently test proposed analytical models against a simulation approach. These model vetting techniques are examples of which of the following vetting phases?
    Accepting/rejecting a model Testing models against simulation\\
    Alpha Capital有限责任公司（Alpha）使用月度模型审查来缓解潜在的模型风险。Alpha的管理者最近接受了使用一个模型来估值25年期公司债券的短期期权，但拒绝了使用同一模型来估值五年期政府债券的短期期权。管理者还经常将提议的分析模型与模拟方法进行测试。这些模型审查技术是以下哪些审查阶段的例子？
    接受/拒绝模型 将模型与模拟进行测试
\end{question}

\begin{choices}
    \item Health check of the model Stress testing\\
    模型健康检查 压力测试
    \item Soundness of a model Stress testing\\
    模型健全性检验 压力测试
    \item Health check of the model Benchmark modeling\\
    模型健康检查 基准建模
    \item Soundness of a model Benchmark modeling\\
    模型健全性检验 基准建模
\end{choices}

\begin{correctanswer}
    D
\end{correctanswer}

\begin{explanation}
    \textbf{A选项错误。}
    \begin{itemize}
        \item 健康检查（Health check）确保模型包含必要的属性和特征，但接受/拒绝模型是健全性检验而非健康检查；压力测试（Stress testing）检查模型在极端条件下的表现，但用模拟方法测试是基准建模而非压力测试。
    \end{itemize}

    \textbf{B选项错误。}
    \begin{itemize}
        \item 模型健全性检验正确，但压力测试不是用模拟方法测试的正确描述；压力测试关注极端情景，而基准建模关注与独立方法的比较验证。
    \end{itemize}

    \textbf{C选项错误。}
    \begin{itemize}
        \item 健康检查不是接受/拒绝模型的正确描述；基准建模正确，但健康检查与健全性检验不同。
    \end{itemize}

    \textbf{D选项正确。}
    \begin{itemize}
        \item 模型健全性检验（Soundness）用于评估模型是否适用于特定资产类型，体现在接受/拒绝模型的决策中；基准建模（Benchmark modeling）用于通过模拟方法等独立方法验证分析模型的准确性。
    \end{itemize}
\end{explanation}

\checklist{模型验证的不同阶段}


\begin{question}
    Which of the following risk aggregation methodologies is characterized by great difficulty in validating parameterization and building a joint distribution?
    \\以下哪种风险加总方法在验证参数化和构建联合分布方面存在极大困难？
\end{question}

\begin{choices}
    \item Copulas\\
    联结函数
    \item Constant diversification\\
    固定分散化
    \item Variance-covariance matrix\\
    方差-协方差矩阵
    \item Full modeling/Simulation\\
    全模型/模拟
\end{choices}

\begin{correctanswer}
    A
\end{correctanswer}

\begin{explanation}
    \textbf{A选项正确。}
    \begin{itemize}
        \item Copula方法在参数验证和构建联合分布方面存在最大困难；Copula方法需要选择Copula类型（如Gaussian Copula、t-Copula、Archimedean Copula等），估计多个参数（如相关系数矩阵、尾部依赖参数、自由度等），这些参数难以从历史数据中准确验证和估计，特别是尾部依赖参数；构建联合分布需要将边际分布与Copula函数结合，过程复杂且容易出错；参数的不确定性也导致验证困难。
    \end{itemize}

    \textbf{B选项错误。}
    \begin{itemize}
        \item 固定分散化方法相对简单，不需要复杂的参数验证；该方法假设固定的分散化系数，不需要估计复杂的参数或构建联合分布，实施和验证相对直接。
    \end{itemize}

    \textbf{C选项错误。}
    \begin{itemize}
        \item 方差-协方差矩阵方法虽然需要估计相关系数，但参数验证相对容易；该方法使用相关系数矩阵，参数估计相对直接，可以通过历史数据计算，联合分布的构建基于线性相关假设，比Copula方法简单。
    \end{itemize}

    \textbf{D选项错误。}
    \begin{itemize}
        \item 全模型/模拟方法虽然需要建模，但参数验证相对直接；该方法通过模拟方法构建联合分布，参数可以通过历史数据或专家判断直接设置，验证过程相对透明和直观，不需要复杂的Copula参数估计。
    \end{itemize}
\end{explanation}

\checklist{风险加总方法的参数验证难度}


\begin{question}
    Which of the following model validation processes is specifically characterized by the limitation that it provides little comfort that the model actually reflects reality?
    \\以下哪种模型验证过程特别具有局限性，即提供的模型反映现实程度的信心很少？
\end{question}

\begin{choices}
    \item Backtesting\\
    回测
    \item Benchmarking\\
    基准测试
    \item Stress testing\\
    压力测试
    \item Qualitative review
    定性审查
\end{choices}

\begin{correctanswer}
    B
\end{correctanswer}

\begin{explanation}
    \textbf{A选项错误。}
    \begin{itemize}
        \item 回测（Backtesting）通过比较模型预测与实际发生的结果来验证模型，能够直接反映模型对现实的拟合程度；回测使用历史数据检验模型预测的准确性，如果模型预测与实际结果一致，说明模型较好地反映了现实；因此回测能提供模型反映现实的验证。
    \end{itemize}

    \textbf{B选项正确。}
    \begin{itemize}
        \item 基准测试（Benchmarking）的一个关键局限性是它几乎不能提供模型是否反映现实的验证；基准测试是将模型结果与另一个模型或方法进行比较，这只是模型间的相对比较，而非与实际数据的比较；如果两个模型都错误，基准测试可能显示它们一致，但实际上都不反映现实；基准测试只能验证模型之间的可比性或参数的一致性，但不能确认模型是否准确反映现实。
    \end{itemize}

    \textbf{C选项错误。}
    \begin{itemize}
        \item 压力测试（Stress testing）虽然不能直接验证历史拟合度，但能够检验模型在极端条件下是否捕捉到现实的极端风险特征；如果模型在压力情景下的表现合理，说明模型能够反映现实的风险特征；因此压力测试也能提供关于模型反映现实的验证。
    \end{itemize}

    \textbf{D选项错误。}
    \begin{itemize}
        \item 定性审查（Qualitative review）通过专家判断评估模型的假设、逻辑和合理性，能够验证模型是否反映现实；虽然不如定量方法客观，但专家审查可以通过评估模型假设的合理性、逻辑的连贯性等来判断模型是否反映现实；因此定性审查也能提供关于模型反映现实的验证。
    \end{itemize}
\end{explanation}

\checklist{模型验证方法的局限性}



\begin{question}
    A bank credit officer, who has reviewed a loan application, has made the following statement: "On a standalone basis, I was not very keen on granting this loan; however, I granted this loan after looking at the overall asset portfolio of the bank." Based on the above statement, which of the following is true?
    \\一个信贷官员审查了一个贷款申请，做出了以下陈述：“在独立基础上，我对这个贷款不太感兴趣；然而，在看了银行的整体资产组合后，我批准了这个贷款。”根据上述陈述，以下哪项是正确的？
\end{question}

\begin{choices}
    \item The correlation of the newly granted loan with the overall portfolio is low and therefore the credit officer was right in granting the loan\\
    新贷款与整体组合的相关性低，因此信贷员批准贷款是正确的
    \item The correlation of the newly granted loan with the overall portfolio is low and therefore the credit officer was wrong in granting the loan\\
    新贷款与整体组合的相关性低，因此信贷员批准贷款是错误的
    \item The correlation of the newly granted loan with the overall portfolio is high and therefore the credit officer was right in granting the loan\\
    新贷款与整体组合的相关性高，因此信贷员批准贷款是正确的
    \item The correlation of the newly granted loan with the overall portfolio is high and therefore the credit officer was wrong in granting the loan\\
    新贷款与整体组合的相关性高，因此信贷员批准贷款是错误的
\end{choices}

\begin{correctanswer}
    A
\end{correctanswer}

\begin{explanation}
    \textbf{A选项正确。}
    \begin{itemize}
        \item 新贷款与整体组合的相关性低，因此信贷员批准贷款是正确的；低相关性意味着分散化收益，虽然单独看贷款风险较高，但低相关性使其加入组合后能够降低组合总风险（通过分散化效应）；贷款的风险由系统性和非系统性风险组成，低相关性意味着非系统性风险可以通过组合分散，因此信贷员的决策合理。
    \end{itemize}

    \textbf{B选项错误。}
    \begin{itemize}
        \item 如果新贷款与组合相关性低，这有利于分散风险，不应拒绝贷款；低相关性提供分散化收益，即使单独看风险较高，也应批准贷款以优化组合风险回报。
    \end{itemize}

    \textbf{C选项错误。}
    \begin{itemize}
        \item 如果新贷款与组合相关性高，这会增加组合风险而非降低；高相关性意味着没有分散化收益，如果单独看风险就高，加入高相关性组合只会增加总风险，因此不应批准贷款。
    \end{itemize}

    \textbf{D选项错误。}
    \begin{itemize}
        \item 虽然高相关性确实不利，但信贷员实际上批准了贷款，说明相关性应该是低的；如果相关性高且单独看风险高，信贷员不应该在看了组合后还批准贷款；这与信贷员的行为矛盾。
    \end{itemize}
\end{explanation}

\checklist{贷款组合风险分散化原则}


\begin{question}
    Which of the following is not a major challenge for banks in reporting measures of conduct and culture?
    \\以下哪项不是银行在报告行为和文化措施方面的主要挑战？
\end{question}

\begin{choices}
    \item Quality data are often not available to enable boards and management to drill down into detail at the business level\\
    高质量数据通常不可用，无法让董事会和管理层深入业务层面的细节
    \item Standardized measures are more difficult to define for smaller regional banks than for global banks\\
    标准化措施对小型区域银行比全球银行更难定义
    \item Limited, asymmetric data make comparisons difficult for identifying trends and relationships with other variables\\
    有限且不对称的数据使得比较困难，难以识别趋势和与其他变量的关系
    \item Management teams at banks are inexperienced in implementing a reporting process for conduct and culture\\
    银行管理团队在实施行为和文化报告流程方面缺乏经验
\end{choices}

\begin{correctanswer}
    B
\end{correctanswer}

\begin{explanation}
    \textbf{A选项是主要挑战。}
    \begin{itemize}
        \item 高质量数据通常不可用，无法让董事会和管理层深入业务层面的细节，这是报告行为和文化指标的主要挑战；行为和文化是定性概念，难以量化，缺乏结构化数据，影响管理层深入分析和业务层面的监控。
    \end{itemize}

    \textbf{B选项不是主要挑战（正确答案）。}
    \begin{itemize}
        \item 标准化指标对小型区域银行比全球银行更难定义不是主要挑战；标准化指标定义困难对所有银行（无论规模大小）都是普遍挑战，而非特定于小型区域银行相对于全球银行；主要挑战在于跨业务线、地区和文化背景的一致性定义，而非规模差异；该表述本身存在问题，因为标准化指标定义的困难与银行规模无直接关系。
    \end{itemize}

    \textbf{C选项是主要挑战。}
    \begin{itemize}
        \item 有限且不对称的数据使得比较困难，难以识别趋势和与其他变量的关系，这是报告行为和文化指标的主要挑战；行为和文化数据通常不完整、不对称（不同来源、不同质量），导致跨时间、跨业务线或跨机构的比较困难，难以识别趋势和建立因果关系。
    \end{itemize}

    \textbf{D选项是主要挑战。}
    \begin{itemize}
        \item 银行管理团队在实施行为和文化报告流程方面缺乏经验，这是主要挑战；行为和文化报告是相对较新的领域，许多银行缺乏成熟的报告框架、流程和经验，需要从零开始建立报告体系和方法。
    \end{itemize}
\end{explanation}

\checklist{银行行为与文化报告的主要挑战}

\begin{question}
    The use of which of the following items is meant more for protecting against risk deterioration?
    \\以下哪项主要用于防范风险恶化？
\end{question}

\begin{choices}
    \item Risk based pricing\\
    风险定价
    \item Management incentives\\
    管理激励
    \item Credit portfolio management\\
    信用组合管理
    \item Customer profitability analysis\\
    客户盈利分析
\end{choices}

\begin{correctanswer}
    C
\end{correctanswer}

\begin{explanation}
    \textbf{A选项错误。}
    \begin{itemize}
        \item 风险定价（Risk-based pricing）主要用于根据风险水平调整价格，以最大化银行利润和风险调整后收益，而非主要用来防范风险恶化；风险定价是定价工具，虽然高风险对应高价格，但主要目的是利润最大化而非风险控制。
    \end{itemize}

    \textbf{B选项错误。}
    \begin{itemize}
        \item 管理激励（Management incentives）主要用于鼓励管理者参与经济资本配置、优化资源配置和提升绩效，而非主要用来防范风险恶化；管理激励是激励工具，主要目的是引导行为，而非直接防范风险恶化。
    \end{itemize}

    \textbf{C选项正确。}
    \begin{itemize}
        \item 信用组合管理（Credit portfolio management）主要用于防范风险恶化；通过组合管理，银行可以监控组合风险水平、识别风险集中度、调整风险敞口、实施分散化策略、设定风险限额、进行风险对冲等，这些措施直接旨在防范组合风险恶化；信用组合管理是主动的风险控制工具，专门用于防范和管理风险恶化。
    \end{itemize}

    \textbf{D选项错误。}
    \begin{itemize}
        \item 客户盈利分析（Customer profitability analysis）主要用于识别不盈利客户、优化客户关系和定价决策，而非主要用来防范风险恶化；这是盈利分析工具，主要目的是利润优化，虽然可能间接影响风险，但不是专门用于防范风险恶化的工具。
    \end{itemize}
\end{explanation}

\checklist{风险恶化防范工具的选择}



\begin{question}
    A risk manager at bank ABC is evaluating approaches the bank can use to measure the total risk in its portfolio. The bank's portfolio contains foreign currency forwards and swaps, credit instruments and derivatives, as well as various types of loans. The manager is studying how the compartmentalized and the unified risk measurement approaches are applied and the differences between them. Which of the following is correct?
    \\银行ABC的风险经理正在评估银行可以使用的测量组合总风险的方法。银行组合包含外汇远期和掉期、信用工具和衍生品，以及各种类型的贷款。经理正在研究分块方法和统一风险测量方法的应用以及它们之间的差异。以下哪项是正确的？
\end{question}

\begin{choices}
    \item If diversification benefits are present in the portfolio, the compartmentalized approach will typically generate a lower estimate of the bank's total risk than the unified approach\\
    如果组合中存在分散化收益，分块方法通常会产生比统一方法更低的银行总风险估计
    \item Interactions between risk categories that result in compounding effects are more appropriately captured by using the compartmentalized approach rather than the unified approach\\
    风险类别间相互作用导致的复合效应更适于使用分块方法而非统一方法捕捉
    \item The unified approach sums the amount of risk separately measured in each risk category into a single estimate of the bank's total risk\\
    统一方法将分别测量的各风险类别风险相加成单一总风险估计
    \item All else held constant, if wrong-way risk is present in the portfolio, the unified approach will typically generate a higher estimate of the bank's total risk than the compartmentalized approach\\
    如果组合中存在错误方向风险，统一方法通常会产生比分块方法更高的银行总风险估计
\end{choices}

\begin{correctanswer}
    D
\end{correctanswer}

\begin{explanation}
   
    \textbf{A选项错误。}
    \begin{itemize}
        \item 如果存在分散化收益，分块方法通常会产生更高的总风险估计而非更低；分块方法简单相加各风险类别，假设完全正相关，忽略分散化收益，会高估总风险；统一方法考虑相关性，能够反映分散化收益，通常产生较低的总风险估计。
    \end{itemize}

    \textbf{B选项错误。}
    \begin{itemize}
        \item 导致复合效应的风险类别间相互作用在统一方法中更好捕捉，而非分块方法；统一方法能够捕捉不同风险类别间的相互作用、相关性增强和复合效应（如信用风险与市场风险的相互作用）；分块方法分别处理各风险类别，可能忽略这些相互作用。
    \end{itemize}

    \textbf{C选项错误。}
    \begin{itemize}
        \item 将分别测量的各风险类别风险相加成单一总风险估计是分块方法的特点，而非统一方法；统一方法会考虑风险类别间的相关性，使用更复杂的聚合方法（如考虑相关性的平方根公式），而非简单相加。
    \end{itemize}

    \textbf{D选项正确。}
    \begin{itemize}
        \item 如果存在错误方向风险（wrong-way risk），统一方法通常会产生更高的总风险估计；错误方向风险是指风险敞口与交易对手信用质量呈负相关的情况（如交易对手违约时敞口增加），统一方法能够捕捉这种风险类别间的负面相互作用和复合效应；而分块方法可能忽略这些相互作用，低估总风险；因此统一方法会产生更高的总风险估计。
    \end{itemize}
\end{explanation}

\checklist{风险度量方法的比较}




\begin{question}
    The CRO of a small bank has recommended that the bank develop an economic capital framework. In a meeting with the board of directors, the CRO lists several potential benefits to the bank from adopting the framework. Which of the following correctly describes a potential benefit of implementing an economic capital framework?
    \\一家小型银行的CRO建议银行发展一个经济资本框架。在与董事会成员的会议上，CRO列出了采用框架对银行潜在的几个好处。以下哪项正确描述了实施经济资本框架的潜在好处？
\end{question}

\begin{choices}
    \item It increases the pricing power of banks by allowing them to charge more competitive interest rates on wholesale loans and retail deposits\\
    它通过允许银行对批发贷款和零售存款收取更具竞争力的利率来增强银行的定价能力
    \item It can be extended down to the level of a single customer in order to assess the customer's profitability on a risk-adjusted basis\\
    它可以扩展到单个客户层面，以评估客户在风险调整基础上的盈利性
    \item It can incentivize senior managers by increasing their compensation as the level of economic capital allocated to their business unit increases\\
    它可以通过增加分配给业务单元的经济资本水平来激励高级管理人员
    \item It allocates capital based on the portfolio's standalone risk rather than its marginal contribution to the bank's total risk\\
    它基于投资组合的独立风险而非其对银行总风险的边际贡献分配资本
\end{choices}

\begin{correctanswer}
    B
\end{correctanswer}

\begin{explanation}

        \textbf{A选项错误。}
        \begin{itemize}
            \item 经济资本框架不直接提高银行的定价能力，也不允许银行通过"更具竞争力的利率"来增加定价权力；经济资本框架有助于风险调整定价（RAROC），但定价能力主要受市场因素、竞争和监管影响；高风险贷款应收取更高利率以补偿风险，而非更低利率；该表述误解了经济资本框架的作用。
        \end{itemize}
    
        \textbf{B选项正确。}
        \begin{itemize}
            \item 经济资本框架可以扩展到单个客户层面，用于评估客户在风险调整基础上的盈利性；这是经济资本框架的重要优势，可以计算单个客户的风险调整资本回报率（RAROC），评估客户关系是否盈利，识别不盈利客户，优化客户关系管理，支持客户层面的定价和资源分配决策。
        \end{itemize}
    
        \textbf{C选项错误。}
        \begin{itemize}
            \item 经济资本增加不应导致高级管理人员薪酬增加；经济资本增加通常意味着风险增加，不应被奖励；薪酬应基于风险调整后的收益（RAROC），而非经济资本本身；高RAROC应获得奖励，而非高经济资本；该表述误解了激励机制的合理设计。
        \end{itemize}
    
        \textbf{D选项错误。}
        \begin{itemize}
            \item 经济资本应基于投资组合对银行总风险的边际贡献，而非仅独立风险；边际贡献反映了新增风险对总风险的增量影响，考虑了分散化效应；仅基于独立风险会高估或低估资本要求，不能正确反映风险的真实成本；该表述误解了经济资本分配的原则。
        \end{itemize}
    \end{explanation}


\checklist{经济资本框架的实施效益}





\section{考点3:大型银行控股公司的资本规划}
\setcounter{question}{0}
\begin{question}
    The Basel Committee recommends that banks use a set of early warning indicators in order to identify emerging risks and potential vulnerabilities in its liquidity position. Which of the following are not early warning indicators of a potential liquidity problem?
    \begin{itemize}
    \item I. Rapid asset growth
    \item II. Negative publicity
    \item III. Credit rating downgrade
    \item IV. Increased asset diversification
\end{itemize}
    巴塞尔委员会建议银行使用一组早期预警指标，以识别新兴风险和潜在的流动性问题。以下哪些不是潜在流动性问题的早期预警指标？
    \begin{itemize}
        \item I. 快速资产增长
        \item II. 负面宣传
        \item III. 信用评级下调
        \item IV. 资产多样化增加
    \end{itemize}
\end{question}

\begin{choices}
    \item II and III
    \item IV only
    \item I and IV
    \item I, II and IV
\end{choices}

\begin{correctanswer}
    B
\end{correctanswer}

\begin{explanation}
    \textbf{指标I分析：是早期预警指标}
    \begin{itemize}
        \item 快速资产增长是流动性问题的早期预警指标；如果资产增长快于稳定的资金来源，可能导致流动性错配；如果资产增长依赖短期或波动性负债融资，在压力时期这些资金来源可能迅速撤离，引发流动性危机。
    \end{itemize}

    \textbf{指标II分析：是早期预警指标}
    \begin{itemize}
        \item 负面宣传是流动性问题的早期预警指标；负面宣传可能导致存款流失、融资渠道受限、市场信心下降、交易对手要求更多抵押品，这些都是流动性危机的征兆。
    \end{itemize}

    \textbf{指标III分析：是早期预警指标}
    \begin{itemize}
        \item 信用评级下调是流动性问题的早期预警指标；信用评级下调可能导致融资成本上升、融资渠道受限、市场信心下降、存款流失、交易对手要求更多抵押品，增加流动性压力。
    \end{itemize}

    \textbf{指标IV分析：不是早期预警指标}
    \begin{itemize}
        \item 资产多样化增加不是流动性问题的早期预警指标；资产多样化通常是有益的，有助于分散风险、降低集中度风险；相反，资产多样化减少或集中度增加才是流动性问题的预警指标；该指标与流动性问题的征兆相反。
    \end{itemize}
\end{explanation}

\checklist{流动性风险的早期预警指标}






\section{考点4:银行压力测试}
\setcounter{question}{0}

\begin{question}
    Which of the following statements regarding stress testing and firm wide risk aggregation is most accurate?
    \\以下关于压力测试和公司层面风险加总的陈述哪项最准确？
\end{question}

\begin{choices}
    \item One potential problem with the risk aggregation process is the underestimation of the diversification effect\\
    风险加总过程中的常见问题是高估而非低估多样化效应
    \item In the context of risk aggregation, stress tests are considered a more reliable alternative than regulatory and economic capital measures\\
    压力测试不应被视为比监管和经济资本计量更可靠的替代方案
    \item The process of stress testing can be performed in a relatively objective manner given that many common risks are reasonably quantifiable\\
    压力测试过程可以相对客观地进行，因为许多常见风险可以合理量化
    \item Using the results of stress testing, those scenarios that will result in minimal losses with a very low likelihood of occurrence may be ignored when considering any adjustments to risk appetite\\
    使用压力测试结果，那些将导致最小损失且发生概率非常低的情景在考虑风险偏好调整时可能被忽略
\end{choices}

\begin{correctanswer}
    D
\end{correctanswer}

\begin{explanation}
    \textbf{A选项错误。}
    \begin{itemize}
        \item 风险加总过程中的常见问题是高估而非低估多样化效应；在压力时期，市场相关性会上升，多样化效果会减弱或消失，风险加总模型往往在正常时期假设的多样化效应在压力时期不再有效，导致风险低估；低估多样化效应的情况较少见。
    \end{itemize}

    \textbf{B选项错误。}
    \begin{itemize}
        \item 压力测试不应被视为比监管和经济资本计量更可靠的替代方案；压力测试是监管资本和经济资本的补充工具，用于评估极端情景下的风险，而不是替代这些资本计量方法；监管资本和经济资本是持续的风险度量，而压力测试是情景分析工具。
    \end{itemize}

    \textbf{C选项错误。}
    \begin{itemize}
        \item 压力测试过程不能以相对客观的方式进行；虽然某些风险可以量化，但压力测试涉及大量主观判断，包括情景设计、假设设定、参数选择、相关性假设等，这些都需要专家判断和管理层决策，难以做到完全客观；压力测试的客观性有限。
    \end{itemize}

    \textbf{D选项正确。}
    \begin{itemize}
        \item 使用压力测试结果，那些将导致最小损失且发生概率非常低的情景在考虑风险偏好调整时可能被忽略；这是合理的风险管理实践，因为风险管理应聚焦主要风险情景和重大损失情景，损失极小且概率极低的情景对风险偏好的影响可以忽略，有助于提高风险管理效率并聚焦关键风险。
    \end{itemize}
\end{explanation}

\checklist{压力测试与风险加总的应用原则}


\begin{question}
    Nordlandia is a country with a developed economy maintaining its own currency, the Nordlandian crown (NLC), and whose most important export is domestically produced oil and natural gas. In a recent stress test of Nordlandia's banking system, several scenarios were considered. Which of the following is most consistent with being part of a coherent scenario?
    \\Nordlandia是一个拥有发达经济体和自己的货币Nordlandian crown (NLC)的国家，其最重要的出口是国内生产的石油和天然气。最近对Nordlandia银行系统的压力测试考虑了几个情景。以下哪项最符合成为连贯情景的一部分？
\end{question}

\begin{choices}
    \item An increase in domestic inflation and appreciation of the NLC\\
    国内通胀上升和NLC升值
    \item A significant increase in crude oil prices and a decrease in the Nordlandian housing price index\\
    原油价格大幅上涨和房价指数下跌
    \item A drop in crude oil prices and appreciation of the NLC\\
    原油价格下跌和NLC升值
    \item A sustained decrease in natural gas prices and a decrease in the Nordlandian stock index\\
    天然气价格持续下跌和股市指数下跌
\end{choices}

\begin{correctanswer}
    D
\end{correctanswer}

\begin{explanation}
    \textbf{A选项错误。}
    \begin{itemize}
        \item 国内通胀上升和NLC升值不一致；通胀上升通常会导致货币购买力下降和货币贬值；虽然作为主要出口石油和天然气的国家，如果出口价格上涨可能带来货币升值，但通胀上升本身通常与货币升值不一致，除非有特殊因素（如出口价格大幅上涨）抵消通胀影响。
    \end{itemize}

    \textbf{B选项错误。}
    \begin{itemize}
        \item 原油价格大幅上涨和房价指数下跌不一致；作为以原油为主要出口的国家，原油价格上涨通常会导致：出口收入增加、贸易顺差扩大、经济繁荣、收入增加、货币升值，进而推动房价上涨；原油价格上涨与房价下跌在经济逻辑上不一致。
    \end{itemize}

    \textbf{C选项错误。}
    \begin{itemize}
        \item 原油价格下跌和NLC升值不一致；作为以原油为主要出口的国家，原油价格下跌通常会导致：出口收入减少、贸易顺差缩小或出现逆差、经济衰退、货币贬值；原油价格下跌与货币升值在经济逻辑上不一致。
    \end{itemize}

    \textbf{D选项正确。}
    \begin{itemize}
        \item 天然气价格持续下跌和股市指数下跌是一致的；作为以天然气为主要出口的国家，天然气价格持续下跌会导致：出口收入减少、贸易顺差缩小、经济衰退、企业盈利下降、投资者信心下降；经济衰退和盈利下降会导致股市下跌；这是一个逻辑连贯的压力测试情景，符合经济因果关系。
    \end{itemize}
\end{explanation}

\checklist{压力测试场景设计的连贯性原则}


\begin{question}
    In a market crash the following are usually true? 
    \begin{itemize}
    \item I. Fixed-income portfolios hedged with short Treasury bonds and futures lose less than those hedged with interest rate swaps given equivalent durations. 
    \item II. Bid–offer spreads widen because of lower liquidity. 
    \item III. The spreads between off-the-run bonds and benchmark issues widen.
    \end{itemize}
    在市场崩盘时，以下哪项通常是正确的？
    \begin{itemize}
        \item I. 用短期国债和期货对冲的固定收益组合损失通常大于用利率互换对冲的组合（给定相同久期）。
        \item II. 买卖价差扩大因为流动性下降。
        \item III. 非主流债券与基准债券之间的利差扩大。
    \end{itemize}
\end{question}

\begin{choices}
    \item I, II, and III
    \item II and III
    \item I and III
    \item None of the above
\end{choices}

\begin{correctanswer}
    B
\end{correctanswer}

\begin{explanation}
    \textbf{陈述I分析：错误}
    \begin{itemize}
        \item 市场崩盘时，用短期国债和期货对冲的固定收益组合损失通常大于用利率互换对冲的组合（给定相同久期）；市场崩盘时发生"flight to quality"（资金逃向高质量资产），国债作为最安全的资产，其价格上涨幅度通常大于利率互换；用国债和期货做空对冲时，国债价格上涨会导致做空头寸损失更大；因此陈述I是错误的。
    \end{itemize}

    \textbf{陈述II分析：正确}
    \begin{itemize}
        \item 市场崩盘时，买卖价差扩大因为流动性下降；市场崩盘时流动性急剧下降，市场参与者更不愿意提供流动性，交易成本上升，买卖价差显著扩大；这是市场崩盘的典型特征。
    \end{itemize}

    \textbf{陈述III分析：正确}
    \begin{itemize}
        \item 非主流债券（off-the-run bonds）与基准债券（benchmark issues）之间的利差扩大；市场崩盘时发生"flight to quality"，投资者偏好高质量、流动性好的基准债券，非主流债券的流动性较差、风险更高，两者之间的信用利差和流动性利差都会扩大；这是市场压力的典型表现。
    \end{itemize}

\end{explanation}

\checklist{市场崩盘时的流动性与利差变化}


\begin{question}
    Piper Hook, a bank examiner, is trying to make sense of stress tests done by one of the banks she examines. The stress tests are multi-factored and complex. The bank is using multiple extreme scenarios to test capital adequacy, making it difficult for Hook to interpret the results. One of the key stress test design challenges that Hook must deal with in her examination of stress tests is
    \\Piper Hook，一个银行检查员，试图理解她检查的一家银行进行的压力测试。压力测试是多因素和复杂的。银行使用多个极端情景来测试资本充足性，使得Hook难以解释结果。Hook在检查压力测试时必须处理的关键压力测试设计挑战是
\end{question}

\begin{choices}
    \item Multiplicity\\
    多重性
    \item Efficiency\\
    效率
    \item Coherence\\
    连贯性
    \item Efficacy\\
    有效性
\end{choices}

\begin{correctanswer}
    C
\end{correctanswer}

\begin{explanation}
    \textbf{A选项错误。}
    \begin{itemize}
        \item 多重性（Multiplicity）虽然指情景数量多，但这不是题目描述的主要挑战；题目强调"难以解释结果"和"多个极端情景"，主要问题不在于情景数量，而在于情景之间的关系是否合理和一致。
    \end{itemize}

    \textbf{B选项错误。}
    \begin{itemize}
        \item 效率（Efficiency）不是压力测试设计的核心挑战；效率关注执行速度和资源利用，而题目描述的重点是多个极端情景的复杂性和结果解释的困难，这更多指向情景设计的合理性问题。
    \end{itemize}

    \textbf{C选项正确。}
    \begin{itemize}
        \item 连贯性（Coherence）是压力测试设计的关键挑战；在多因素、复杂的压力测试中，确保各个因素之间的关系在经济逻辑上一致、合理且可能发生是核心难点；例如，如果同时设定利率上升、货币升值和经济衰退，需要确保这些因素在经济上是一致的，而非相互矛盾；缺乏连贯性的情景会导致结果难以解释，因为情景本身可能不合理或不可能发生。
    \end{itemize}

    \textbf{D选项错误。}
    \begin{itemize}
        \item 有效性（Efficacy）虽然重要，但不是题目描述的主要设计难点；题目强调"难以解释结果"，这主要指向情景的连贯性问题，即情景是否合理、逻辑一致，而非情景是否有效达到测试目标。
    \end{itemize}
\end{explanation}

\checklist{压力测试设计中的连贯性原则}


\begin{question}
    Sarah Chen, a senior analyst at the Federal Reserve Bank, is presenting research on stress testing methodologies to a group of banking supervisors. She explains that macro-variable stress testing can be misleading for certain banks due to regional variations in economic indicators. She cites the example of significant differences in unemployment rates across European countries following the 2010-2012 sovereign debt crisis. Which type of loan did Chen most likely identify as having the strongest correlation with unemployment rates?
    \\Sarah Chen，美联储银行的高级分析师，正在向一组银行监管者介绍压力测试方法。她解释说，由于经济指标的地区差异，宏观变量压力测试可能对某些银行产生误导。她引用了2010-2012年主权债务危机期间欧洲各国失业率显著差异的例子。Chen最有可能识别出与失业率相关性最强的贷款类型是什么？
\end{question}

\begin{choices}
    \item Residential mortgage loans\\
    住房抵押贷款
    \item Credit card loans\\
    信用卡贷款
    \item Commercial property loans\\
    商业地产贷款
    \item Corporate term loans\\
    企业定期贷款
\end{choices}

\begin{correctanswer}
    B
\end{correctanswer}

\begin{explanation}
    \textbf{A选项错误。}
    \begin{itemize}
        \item 住房抵押贷款虽然受失业率影响，但与失业率的相关性不如信用卡贷款强；住房抵押贷款通常有抵押品保护，违约率对失业率变化的敏感度相对较低；即使失业导致还款困难，借款人可能仍能通过出售房产偿还贷款，违约率上升相对缓慢。
    \end{itemize}

    \textbf{B选项正确。}
    \begin{itemize}
        \item 信用卡贷款与失业率的相关性最强；信用卡贷款是无担保贷款，没有抵押品保护，违约率对失业率变化最敏感；失业直接影响消费者的可支配收入和还款能力，失业率上升时信用卡违约率会显著上升；在2010-2012年主权债务危机期间，希腊失业率达27\%而德国仅5\%，信用卡违约率在两国间存在显著差异，这解释了为什么宏观变量压力测试可能因地区差异而产生误导。
    \end{itemize}

    \textbf{C选项错误。}
    \begin{itemize}
        \item 商业地产贷款虽然受经济周期影响，但与失业率的相关性不如信用卡贷款强；商业地产贷款违约主要受商业活动、租金收入、市场供需影响，与失业率存在间接关系，但相关性相对较弱；商业地产贷款违约通常滞后于失业率变化，且受其他因素（如商业地产市场周期）影响更大。
    \end{itemize}

    \textbf{D选项错误。}
    \begin{itemize}
        \item 企业定期贷款与失业率的相关性较弱；企业贷款违约主要受行业因素、企业财务状况、现金流、盈利能力影响，与失业率存在间接关系（失业率影响整体经济，进而影响企业），但相关性不如消费者贷款（如信用卡）直接和强烈；企业可以通过多种方式应对经济衰退，违约率对失业率变化的敏感度较低。
    \end{itemize}
\end{explanation}

\checklist{压力测试中不同贷款类型对失业率的敏感度}




\end{document}